\documentclass{article}
\usepackage[utf8]{inputenc}
\usepackage[T1]{fontenc}
\usepackage{array}
\usepackage{calc}
\usepackage{booktabs}
\usepackage{geometry}
\usepackage{graphicx}
\usepackage{longtable}
\usepackage{soul}
\usepackage{hyperref}
\usepackage{amssymb}
\usepackage{unicode-math}
\unimathsetup{math-style=TeX}

\geometry{a4paper}
\title{Using Trie for data representation}
\author{Maksym Voroniy}

\subsection{What Can Be a Trie Key?}\label{what-can-be-a-trie-key}

The final question addressed within the scope of this article concerns
the nature of Trie keys. While the answer "strings" is the most obvious
one, real-world systems - such as key-value databases - often rely on
integers, floating-point numbers, or other structured data types.

A separate chapter will therefore discuss how numeric values can be
transformed and represented as keys in Trie-based structures, enabling
their use for efficient indexing beyond textual data.

